\section{Discussion}
\textit{Interpreting transfer asymmetry.} The FIT partitions show descriptive source differences: HP benign URLs are short (median \res{medLenBBenign} characters) and mostly HTTPS (\res{httpsBBenign}\%), whereas HP phishing URLs are longer (\res{medLenBPhish}) with more dots and subdomains (medians \res{medDotsBPhish} and \res{medSubBPhish}). In EBBU, phishing URLs have medians of \res{medSubAPhish} subdomains and \res{medDotsAPhish} dots and are rarely HTTPS (\res{httpsAPhish}\%), but neither are its benign URLs (\res{httpsABenign}\% HTTPS). The RQ5 diagnostics turn this description into a mechanism. Some cues are shared in direction (HTTPS, length, entropy), others reverse (host-label and subdomain counts, query presence). EBBU-trained models meet HP data in which the main errors are HTTPS phishing, a population that barely exists in EBBU. HP-trained models meet EBBU data whose benign URLs lack the \texttt{www} label and whose phishing URLs carry few subdomains, contrary to what HP taught them. The slice error rates are consistent with this account; they remain correlational and do not prove that a single feature caused an individual error. Rashid~\textit{et al.}~\cite{rashid2024} already discuss feature-distribution differences, so identifying source shift is not our novelty. The two datasets differ in collection process, period and prevalence; our transfer drop describes this benchmark pair rather than isolating a causal domain-shift effect.

\textit{What does stability buy?} STABLE-PI prioritises features with more positive top-$q$ fold votes; it does not require every selected feature to be stable in all folds and does not guarantee invariance across datasets. It outperformed mean ranking at the same small budget in HP$\rightarrow$EBBU, without a clear advantage in the other direction. The matched-$k$ design separates this ranking comparison from the number of retained features. ALL nevertheless remained the stronger cross-dataset reference. Our empirical insight is therefore conditional: source-domain stability can mitigate the loss from a compact subset in this benchmark, but it does not establish that source-selected compression preserves transfer or that the rule is generally safer. RQ5 shows why. Stability is measured across folds of \emph{one} collection, so it rewards features that are consistently useful there, including features whose class association reverses elsewhere. ALL keeps more redundant cues and so has more ways to recover when one cue fails. Stability across \emph{sources}, rather than across folds, is what would be needed, and that requires a second labelled source (Section~\ref{sec:future}).

\textit{What the two-paper extension establishes.} We connect the classifier/representativeness questions in~\cite{sanchez2022} with the transfer problem in~\cite{rashid2024} through a controlled source-only evaluation. Both primary papers satisfy the five-year recency limit; original EBBU provenance is retained in the repository. The contribution lies in the auditable protocol and the ranking--budget--cost comparison. We have not run either original system under identical conditions, so we cannot attribute improvements over their published results to our approach. The two primary papers satisfy the revised project count and recency conditions; they do not establish universal novelty, and relevant adjacent work remains acknowledged.

\textit{Did we get what we wanted?} Partly. The hypothesis that stability-based selection transfers better than mean importance at equal budgets was supported in one direction and neutral in the other. The broader hope that a compact, cheaper model could keep cross-dataset accuracy was not supported: no subset matched ALL, and the latency gains are modest. The study did deliver a clear negative result and a mechanism, which matters for practice: a lexical detector that looks strong on one collection can depend on period-specific cues such as HTTPS.

\textit{What we would do differently.} We would (i) include LogReg and the one-feature rule as pre-registered baselines rather than post-freeze diagnostics; (ii) plan class-conditional feature analysis of both sources \emph{before} defining the selectors, which would have motivated a cross-source criterion; (iii) budget for a third source from the start, so that selection evidence and evaluation come from different collections; and (iv) separate threshold transfer from ranking transfer in the primary outcome, by reporting AP-based contrasts next to F1.

\textit{Uncertainty.} The bootstrap intervals resample target domains for fixed fitted models; they do not include resampling or split-induced training variability. Table~\ref{tab:contrasts} shows that per-seed differences vary, so the size of the advantage is less stable than its observed sign. Three estimator seeds on a single split do not establish a population-level training effect, and we make no family-wise significance claim. Additional independently collected holdout data and repeated training splits would strengthen robustness claims, but must be planned under a new protocol before opening a new holdout.

\textit{Cost.} Lexical models are already cheap; reducing 37 to 8--16 features trims extraction work but the shared parsing step and the per-call overhead of the model remain, so savings are modest relative to the accuracy given up on the hard direction. ``Lightweight'' is therefore relative to the ALL model on the reported hardware.

\textit{Threats to validity.} (i) Both datasets are older public collections; results do not establish contemporary or zero-day robustness. (ii) The symmetric purge removes shared popular domains and changes the evaluation population. (iii) Both selectors share an RF proxy, which may favour RF. (iv) Correlated lexical features can depress PI. (v) Hyper-parameters were fixed, not optimised. (vi) HP encoding is crossvalidated against the primary paper's class counts in all original splits; a verbatim numeric codebook is not supplied. The complete user-supplied article PDF was reviewed, including its data, alignment algorithm and limitations. (vii) RQ5 is exploratory, post-freeze and computed on VAL of a single split; source-VAL F1 is optimistic because it also sets thresholds. (viii) The freeze is a procedural control, signed by an agent on the author's instruction. (ix) The original SHAP outputs failed additivity and the legacy analysis accessed TEST labels outside the evaluator. Both findings remain in the v1 audit trail. This version uses a separately hashed, label-free exploratory correction on saved feature matrices, with strict additivity and classifier-equivalence checks. Correction after observing v1 outputs does not create a new independent holdout or establish causal explanation fidelity. Full-text review of the primary papers supports this methodological distinction without establishing global novelty.
